% Part: first-order-logic
% Chapter: natural-deduction
% Section: proving-things-quant

\documentclass[../../../include/open-logic-section]{subfiles}

\begin{document}

\olfileid{fol}{ntd}{prq}

\olsection{له سورونو سره \usetoken{P}{derivation}}

\begin{ex}
له سورونو سره د کار پر مهال بايد ډاډ واخلو چې د ځانګړي متغير شرط
نۀ ماتوو، او کله کله د دې دپاره د ځينو استنتاجونو د عملي کولو ترتيب
بدلول را بدلول غواړي۔ په عمومي توګه ګټه کوي چې لومړے د هغو قاعدو
د حل هڅه وکړو چې د ځانګړي متغير شرط ورباندې لګېږي (په بشپړ ثبوت
کښې به دوئ لاندې راځي)۔

راځئ وګورو چې د !!{formula}~$\lexists[x][\lnot !A(x)] \lif \lnot
\lforall[x][!A(x)]$ يو !!a{derivation} څنګه ورکوو۔ د معمول په شان
داسې پيل کوو:
\begin{prooftree}
\AxiomC{}
\UnaryInfC{$\lexists[x][\lnot !A(x)]\lif \lnot \lforall[x][!A(x)]$}
\end{prooftree}
لومړے ليکو چې د~\Intro{\lif} د قاعدې په وسيله د وروستي پړاو د
توجيه کولو دپاره څه پکار دي۔
\begin{prooftree}
\AxiomC{$\Discharge{\lexists[x][\lnot !A(x)]}{1}$}
\DeduceC{$\lnot \lforall[x][!A(x)]$}
\DischargeRule{\Intro{\lif}}{1}
\UnaryInfC{$\lexists[x][\lnot !A(x)]\lif \lnot \lforall[x][!A(x)]$}
\end{prooftree}
ځکه پر~$\lnot \lforall[x][!A(x)]$ د کارولو دپاره څرګنده قاعده نۀ
شته، !!{derivation} داسې برابروو چې د~\Elim{\lexists} قاعده
وکارولے شو۔ دلته بايد د ځانګړي متغير شرط ته پام وکړو او داسې ثابت
وټاکو چې په~$\lexists[x][!A(x)]$ يا د هغې په هرې اړوندې فرضيې کښې
نۀ راځي۔ (خو ځکه چې هېڅ !!{constant}s نۀ دي راغلي، هر انتخاب سم
دے۔)
\begin{prooftree}
\AxiomC{$\Discharge{\lexists[x][\lnot !A(x)]}{1}$}
\AxiomC{$\Discharge{\lnot !A(a)}{2}$}
\DeduceC{$\lnot \lforall[x][!A(x)]$}
\DischargeRule{\Elim{\lexists}}{2}
\BinaryInfC{$\lnot \lforall[x][!A(x)]$}
\DischargeRule{\Intro{\lif}}{1}
\UnaryInfC{$\lexists[x][\lnot !A(x)] \lif \lnot \lforall[x][!A(x)]$}
\end{prooftree}
د~$\lnot \lforall[x][!A(x)]$ د راايستلو دپاره د~\Intro{\lnot}
قاعده کارول غواړو: دا غواړي چې يو تناقض راوباسو، او کېدے شي
$\lforall[x][!A(x)]$ د يوې اضافي فرضيې په توګه وکاروو۔ البته په دې
تناقض کښې هغه فرضيه~$\lnot !A(a)$ هم راتلے شي چې د
\Elim{\lexists} استنتاج به ئې وباسي۔ داسې ئې برابروو:
\begin{prooftree}
\AxiomC{$\Discharge{\lexists[x][\lnot !A(x)]}{1}$}
\AxiomC{$\Discharge{\lnot !A(a)}{2}, \Discharge{\lforall[x][!A(x)]}{3}$}
\DeduceC{$\lfalse$}
\DischargeRule{\Intro{\lnot}}{3}
\UnaryInfC{$\lnot \lforall[x][!A(x)]$}
\DischargeRule{\Elim{\lexists}}{2}
\BinaryInfC{$\lnot \lforall[x][!A(x)]$}
\DischargeRule{\Intro{\lif}}{1}
\UnaryInfC{$\lexists[x][\lnot !A(x)]\lif \lnot \lforall[x][!A(x)]$}
\end{prooftree}
داسې ښکاري چې تناقض ته نژدې يو۔ تر ټولو اسانه قاعده
\Elim{\lforall} ده چې د ځانګړي متغير هېڅ شرط نۀ لري۔ ځکه د عمومي
سور لرونکي~$x$ پر ځاے هره تړلې اصطلاح کارولے شو، ښه ده چې هماغه
~$a$ وکاروو څو تناقض ته ورسېږو۔
\begin{prooftree}
\AxiomC{$\Discharge{\lexists[x][\lnot !A(x)]}{1}$}
\AxiomC{$\Discharge{\lnot !A(a)}{2}$}
\AxiomC{$\Discharge{\lforall[x][!A(x)]}{3}$}
\RightLabel{\Elim{\lforall}}
\UnaryInfC{$!A(a)$}
\RightLabel{\Elim{\lnot}}
\BinaryInfC{$\lfalse$}
\DischargeRule{\Intro{\lnot}}{3}
\UnaryInfC{$\lnot \lforall[x][!A(x)]$}
\DischargeRule{\Elim{\lexists}}{2}
\BinaryInfC{$\lnot \lforall[x][!A(x)]$}
\DischargeRule{\Intro{\lif}}{1}
\UnaryInfC{$\lexists[x][\lnot !A(x)]\lif \lnot \lforall[x][!A(x)]$}
\end{prooftree}

په ځانګړي ډول له سورونو سره د کار پر مهال اوس دا بيا کتل مهم دي
چې د ځانګړي متغير شرط نۀ دے مات شوے۔ دلته دغه شرط يوازې پر
\Elim{\exists} لګېدۀ، او ځانګړے متغير~$a$ په هغو فرضيو کښې نۀ
راځي چې استنتاج ورپورې تړلے دے؛ نو دا سم !!{derivation} دے۔
\end{ex}

\begin{ex}
کله کله له نورو !!{formula}s څخه !!a{formula} راايستلے شو۔ په دغو
حالتونو کښې کېدے شي نۀ ايستل شوې فرضيې ولرو۔ د خپلې وروستۍ موخې
تر څنګ د فرضيو حساب ساتل هم مهم دي۔

راځئ وګورو چې له فرضيو~$\lexists[x][(!A(x) \land !B(x))]$ او
$\lforall[x][(!B(x) \lif !C(x,b))]$ څخه د !!{formula}
$\lexists[x][!C(x,b)]$ يو !!a{derivation} څنګه ورکوو۔ د معمول په
شان پايله لاندې ليکو۔
\begin{prooftree}
\AxiomC{}
\UnaryInfC{$\lexists[x][!C(x,b)]$}
\end{prooftree}

د کار دپاره دوه مقدمې لرو۔ د لومړۍ د کارولو، يعنې له
$\lexists[x][(!A(x) \land !B(x))]$ څخه د
$\lexists[x][!C(x, b)]$ د !!a{derivation} موندلو دپاره به د
\Elim{\lexists} قاعده کاروو۔ ځکه د ځانګړي متغير شرط لري، دا قاعده
لومړۍ کاروو۔ لاندې لاس ته راځي:
\begin{prooftree}
\AxiomC{$\lexists[x][(!A(x) \land !B(x))]$}
\AxiomC{$\Discharge{!A(a) \land !B(a)}{1}$}
\DeduceC{$\lexists[x][!C(x,b)]$}
\DischargeRule{\Elim{\lexists}}{1}
\BinaryInfC{$\lexists[x][!C(x,b)]$}
\end{prooftree}
هغه دوه فرضيې چې ورسره کار کوو، $!B$ شريکوي۔ ښايي اوس د
\Elim{\land} په کارولو د~$!B(a)$ بېلول ګټور وي۔
\begin{prooftree}
\AxiomC{$\lexists[x][(!A(x) \land !B(x)])$}
\AxiomC{$\Discharge{!A(a) 
\land !B(a)}{1}$}
\RightLabel{\Elim{\land}}
\UnaryInfC{$!B(a)$}
\DeduceC{$\lexists[x][!C(x,b)]$}
\DischargeRule{\Elim{\lexists}}{1}
\BinaryInfC{$\lexists[x][!C(x,b)]$}
\end{prooftree}

د کار دپاره دويمه فرضيه~$\lforall[x][(!B(x) \lif !C(x,b))]$ ده۔
ځکه د ځانګړي متغير شرط نۀ لري، د~\Elim{\lforall} په کارولو د
!!{constant}~$a$ په وسيله $x$ بېلګه کولے او~$!B(a) \lif !C(a, b)$
اخستلے شو۔ اوس $!B(a) \lif !C(a,b)$ او~$!B(a)$ دواړه لرو۔ بل ګام
بايد د~\Elim{\lif} نېغه کارونه وي۔
\begin{prooftree}
\AxiomC{$\lexists[x][(!A(x) \land !B(x))]$}
\AxiomC{$\lforall[x][(!B(x) \lif !C(x,b))]$}
\RightLabel{\Elim{\lforall}}
\UnaryInfC{$!B(a) \lif !C(a,b)$}
\AxiomC{$\Discharge{!A(a) 
\land !B(a)}{1}$}
\RightLabel{\Elim{\land}}
\UnaryInfC{$!B(a)$}
\RightLabel{\Elim{\lif}}
\BinaryInfC{$!C(a,b)$}
\DeduceC{$\lexists[x][!C(x,b)]$}
\DischargeRule{\Elim{\lexists}}{1}
\insertBetweenHyps{\hspace{-5em}}
\BinaryInfC{$\lexists[x][!C(x,b)]$}
\end{prooftree}
موخې ته ډېر نژدې يو!{} د~\Intro{\lexists} يوه کارونه موخې ته مو
رسوي۔
\begin{prooftree}
\AxiomC{$\lexists[x][(!A(x) \land !B(x))]$}
\AxiomC{$\lforall[x][(!B(x) \lif !C(x,b))]$}
\RightLabel{\Elim{\lforall}}
\UnaryInfC{$!B(a) \lif !C(a,b)$}
\AxiomC{$\Discharge{!A(a) 
\land !B(a)}{1}$}
\RightLabel{\Elim{\land}}
\UnaryInfC{$!B(a)$}
\RightLabel{\Elim{\lif}}
\BinaryInfC{$!C(a,b)$}
\RightLabel{\Intro{\lexists}}
\UnaryInfC{$\lexists[x][!C(x,b)]$}
\DischargeRule{\Elim{\lexists}}{1}
\insertBetweenHyps{\hspace{-5em}}
\BinaryInfC{$\lexists[x][!C(x,b)]$}
\end{prooftree}
ځکه په هر پړاو کښې مو ډاډ اخستے چې د ځانګړي متغير شرطونه نۀ دي
مات شوي، باور کولے شو چې دا سم !!{derivation} دے۔
\end{ex}

\begin{ex}
له فرضيو~$\lforall[x][!A(x)] \lif \lexists[y][!B(y)]$ او
$\lnot\lexists[y][!B(y)]$ څخه د !!{formula}
$\lnot\lforall[x][!A(x)]$ يو !!a{derivation} ورکړئ۔ د معمول په شان
موخې !!{formula} لاندې ليکو۔
\begin{prooftree}
\AxiomC{}
\UnaryInfC{$\lnot\lforall[x][!A(x)]$}
\end{prooftree}
د !!{derivation} وروستۍ کرښه نفي ده، نو \Intro{\lnot} کارول غواړو۔
دا غواړي چې معلومه کړو تناقض څنګه !!{derive} کړو۔
\begin{prooftree}
\AxiomC{$\Discharge{\lforall[x][!A(x)]}{1}$}
\DeduceC{$\lfalse$}
\DischargeRule{\Intro{\lnot}}{1}
\UnaryInfC{$\lnot\lforall[x][!A(x)]$}
\end{prooftree}
تر دې ځايه ښه يو۔ \Elim{\lforall} کارولے شو، خو څرګنده نۀ ده چې
موخې ته به مو ورسوي۔ پر ځاے ئې خپله يوه فرضيه کاروو۔
$\lforall[x][!A(x)] \lif \lexists[y][!B(y)]$ له
$\lforall[x][!A(x)]$ سره د~\Elim{\lif} قاعدې کارول راکوي۔
\begin{prooftree}
\AxiomC{$\lforall[x][!A(x)] \lif \lexists[y][!B(y)]$}
\AxiomC{$\Discharge{\lforall[x][!A(x)]}{1}$}
\RightLabel{\Elim{\lif}}
\BinaryInfC{$\lexists[y][!B(y)]$}
\DeduceC{$\lfalse$}
\DischargeRule{\Intro{\lnot}}{1}
\UnaryInfC{$\lnot\lforall[x][!A(x)]$}
\end{prooftree}
اوس د کار دپاره يوه وروستۍ فرضيه لرو، او ښکاري چې د~\Elim{\lnot}
په کارولو تناقض ته مو رسولے شي۔
\begin{prooftree}
\AxiomC{$\lnot\lexists[y][!B(y)]$}
\AxiomC{$\lforall[x][!A(x)] \lif \lexists[y][!B(y)]$}
\AxiomC{$\Discharge{\lforall[x][!A(x)]}{1}$}
\RightLabel{\Elim{\lif}}
\BinaryInfC{$\lexists[y][!B(y)]$}
\RightLabel{\Elim{\lnot}}
\BinaryInfC{$\lfalse$}
\DischargeRule{\Intro{\lnot}}{1}
\UnaryInfC{$\lnot\lforall[x][!A(x)]$}
\end{prooftree}
\end{ex}

\begin{prob}
داسې !!{derivation}s ورکړئ چې لاندې وښيي:
\begin{enumerate}
\item $\Proves (\lforall[x][!A(x)] \land \lforall[y][!B(y)]) \lif
\lforall[z][(!A(z) \land !B(z))]$.
\item $\Proves (\lexists[x][!A(x)] \lor \lexists[y][!B(y)]) \lif
\lexists[z][(!A(z) \lor !B(z))]$.
\item $\lforall[x][(!A(x) \lif !B)] \Proves \lexists[y][!A(y)] \lif !B$.
\item $\lforall[x][\lnot !A(x)] \Proves \lnot\lexists[x][!A(x)]$.
\item $\Proves \lnot\lexists[x][!A(x)] \lif \lforall[x][\lnot !A(x)]$.
\item $\Proves \lnot\lexists[x][\lforall[y][((!A(x,y) \lif \lnot
!A(y,y)) \land (\lnot !A(y,y) \lif !A(x,y)))]]$.
\end{enumerate}
\end{prob}

\begin{prob}
داسې !!{derivation}s ورکړئ چې لاندې وښيي:
\begin{enumerate}
\item $\Proves \lnot\lforall[x][!A(x)] \lif \lexists[x][\lnot!A(x)]$.
\item $(\lforall[x][!A(x)] \lif !B) \Proves \lexists[y][(!A(y) \lif !B)]$.
\item $\Proves \lexists[x][(!A(x) \lif \lforall[y][!A(y)])]$.
\end{enumerate}
(دا ټول د~$\FalseCl$ قاعده غواړي۔)
\end{prob}

\end{document}
